\section{Conclusion}
\label{sec:conclusion}

Safety alignment in VLMs is evaluated on clean images, but deployed on corrupted ones. We showed that this gap is not hypothetical: ordinary, non-adversarial corruptions --- blur, snow, compression-level degradations --- systematically weaken the safety behavior of every model we tested, while leaving utility largely intact, and the direction of the effect is governed by a simple principle: corruption hurts safety when the image carries the \emph{evidence} of harm, and helps when the image carries the \emph{attack}. Existing safety tuning partially re-breaks under corruption; corruption-aware training heals the reasoning but not yet the final answer; and a generic safety prompt is a surprisingly strong, free mitigation --- while the corruption-aware prompt we designed on top of it fails to transfer to either held-out model, a cautionary result about how readily a single-model defense can look general. The practical upshot for the community is direct: safety benchmarks and safety training pipelines should treat image corruption the way robustness benchmarks have for years --- as a first-class evaluation axis, not an afterthought.
